% ---------------------------------------------------------------- Replication results (P1)
\subsection{Replication of the option benchmarks}
\label{sec:replication}

\paragraph{Set-up.}
We ran three Markovian problems for which the exact price is known.
\emph{E1}: the Bermudan put of Section~2 ($S_0=36$, $K=40$, $r=0.06$, $\sigma=0.2$, $T=1$,
50 exercise dates; exact price $4.478$ from the tree and the PDE solver).
\emph{E2}: the paper's Bermudan max-call on $d=20$ and $50$ independent assets
($s_0=100$, 9 exercise dates, Table~1 of~\cite{ospg2023}); the reference is the midpoint of
the lower and upper bounds of Becker et al.~\cite{becker2019}, $[51.571, 51.765]$ and
$[69.582, 69.889]$.
\emph{E3}: the paper's American call on the geometric average of $d=7$ and $20$ correlated
assets (100 exercise dates, Appendix~D.3); the geometric average of correlated geometric
Brownian motions is itself one, so the exact price follows from a one-dimensional tree.
Each problem was simulated with 100{,}000 paths (40{,}000 for E2), split at random into
50\% test and 50\% training, of which a fifth is held out for validation, and repeated for
three seeds. Every method sees the same features and is scored on the same test paths.

\begin{table}[ht]
\centering\small
\caption{Out-of-sample values, mean over three seeds (standard deviation across seeds).
The standard error of one test set is about $0.013$ (put), $0.16$ (max-call) and $0.065$
(geometric call). DOS was not run on the 100-date geometric call.}
\label{tab:replication}
% generated by experiments/summarize.py: method & mean (sd) per problem
\begin{tabular}{lrrrrr}
\toprule
Method & Put & Max-call $d{=}20$ & Max-call $d{=}50$ & Geo-call $d{=}7$ & Geo-call $d{=}20$ \\
\midrule
LSM & 4.458 (0.004) & 51.69 (0.20) & 69.83 (0.32) & 10.169 (0.012) & 9.899 (0.011) \\
DOS & 4.465 (0.007) & 51.69 (0.16) & 69.72 (0.29) & -- & -- \\
FQI-MLP & 4.467 (0.011) & 51.68 (0.12) & 69.67 (0.21) & 10.145 (0.068) & 9.874 (0.056) \\
FQI-GRU & -- & 51.73 (0.14) & 69.81 (0.37) & -- & -- \\
OSPG-MLP & 4.473 (0.008) & 51.80 (0.20) & 69.93 (0.28) & 10.255 (0.046) & 9.989 (0.017) \\
OSPG-GRU & 4.465 (0.015) & 51.72 (0.21) & 69.87 (0.25) & 10.237 (0.065) & 9.941 (0.030) \\
\midrule
Reference & 4.478 & 51.67 & 69.74 & 10.260 & 10.033 \\
Paper, DNN-OSPG & -- & 51.02 & 69.06 & 10.252 & 10.010 \\
\bottomrule
\end{tabular}

\end{table}

\paragraph{Every method recovers the price; OSPG-MLP comes closest.}
Table~\ref{tab:replication} shows all methods within $1.6\%$ of the exact price. The
memoryless OSPG policy is the most accurate on the three problems where the differences
are larger than the noise: $-0.12\%$ on the put, and $-0.05\%$ and $-0.44\%$ on the
geometric call, against $-0.9\%$ to $-1.6\%$ for Longstaff--Schwartz and fitted
Q-iteration. This is the paper's ranking (OSPG ahead of FQI and LS), and our OSPG errors are
of the same size as the paper's ($0.07\%$ and $0.23\%$), although on the geometric call we
trained on about thirty times fewer paths. On the max-call every method is within one standard error of the
reference. OSPG-MLP's $51.80$ and $69.93$ lie slightly above Becker et al.'s upper bounds,
but a mean of three test sets has a standard error of about $0.09$, so neither is
significant.

\paragraph{Paired comparisons.}
Because all methods are evaluated on the same test paths, we compare them path by path:
the noise common to both methods cancels and much smaller differences become visible
(\texttt{results/paired.md}). OSPG-GRU is ahead of Longstaff--Schwartz and FQI-MLP on both
geometric calls by $0.04$ to $0.09$ ($z$ between $2.6$ and $4.6$), and of DOS and FQI-MLP on
the 50-asset max-call by $0.15$ and $0.20$ ($z=3.9$ and $5.4$). It is, however, behind
OSPG-MLP on the put ($0.008$, $z=3.0$), on both max-calls ($0.086$ and $0.064$, $z=2.9$ and
$2.2$) and on the 20-asset geometric call ($0.049$, $z=3.7$).

\paragraph{Memory has a small cost when there is nothing to remember.}
In these Markovian problems the history carries no information beyond today's state,
so the recurrent policy has to learn to ignore it, and it does so imperfectly: it loses
$0.1\%$ to $0.5\%$ against the memoryless policy and trains 1.3 to 4 times more slowly.
Part of the gap may come from batch normalisation, which the MLP uses (as in the paper)
and the GRU does not. The loss is small next to the margin by which both OSPG policies beat
the backward-induction methods on the harder problems, so the recurrent policy remains a
reasonable default when it is not known whether the problem has memory.

\paragraph{The paper's own max-call values.}
The paper reports $50.86$--$51.02$ for $d=20$ and $68.87$--$69.06$ for $d=50$, about $0.8$
below Becker et al.'s lower bounds and below every method here, including our
implementations of the paper's baselines. We therefore compare with Becker et al.\ rather than with the
paper's numbers; the paper's conclusions about the max-call (all methods within a few
tenths of each other) are unaffected.

% Note for P3 (rough Bergomi, stand-in simulator, results/paired.md): both OSPG policies
% beat LSM, DOS and FQI-GRU by 0.035-0.08 (z = 2.3-4.5), but OSPG-GRU and OSPG-MLP are equal
% (difference 0.004, z = 0.3). On the stand-in, the gain comes from avoiding backward
% induction, not from reading the history. Re-check with your simulator and an H sweep.
